\Needspace{20\baselineskip}
\section{计算规模与端到端性能}
\label{l12:sec:performance}

\subsection{矩阵形状如何决定运算量}

浮点运算次数（floating-point operations，FLOP）的计数需要先固定约定。把一次实数乘法和一次实数加法各计为一次运算，则一次融合乘加（fused multiply-add，FMA）计为 $2$ FLOP。对于 $(m\times k)(k\times n)$ 的稠密矩阵乘法，主导运算量按 $2mkn$ 计算。下面只统计主要矩阵乘法，暂不计偏置、归一化、激活和采样中的低阶或特殊函数开销。

在标准多头结构 $hd=D$ 下，所有头的三个输入投影合计为 $6ND^2$ FLOP，输出投影为 $2ND^2$ FLOP。单头 $QK^{\mathsf T}$ 需要约 $2N^2d$ FLOP，所有头合计 $2N^2D$；$PV$ 同样需要 $2N^2D$。前馈网络的两次乘法合计 $4NDD_{\mathrm{ff}}$。由此得到表~\ref{l12:tab:flops}。

\begin{table}[htbp]
\centering
\caption{单序列、单层的主要矩阵乘法开销。预填充列采用完整稠密矩阵乘法的计数，解码列的可见长度为 $N$。}
\label{l12:tab:flops}
\begin{tabularx}{\linewidth}{@{}Y L{4.3cm}L{4.3cm}@{}}
\toprule
计算部分 & 长度 $N$ 的预填充 & 一个位置的增量解码\\
\midrule
全部查询、键和值投影 & $6ND^2$ & $6D^2$\\
分数矩阵 $QK^{\mathsf T}$ & $2N^2D$ & $2ND$\\
值聚合 $PV$ & $2N^2D$ & $2ND$\\
多头输出投影 & $2ND^2$ & $2D^2$\\
前馈两次投影 & $4NDD_{\mathrm{ff}}$ & $4DD_{\mathrm{ff}}$\\
\bottomrule
\end{tabularx}
\end{table}

因此，一层预填充的主要算量为 $8ND^2+4N^2D+4NDD_{\mathrm{ff}}$。取 $D_{\mathrm{ff}}=4D$，可写成
\begin{equation}
 F_{\mathrm{prefill}}\approx24ND^2+4N^2D,\qquad
 F_{\mathrm{decode}}\approx24D^2+4ND.
 \label{l12:eq:layer-flops}
\end{equation}
对包含 $L$ 层的网络，层内部分再乘 $L$。最终词表投影需要另外计入：若计算全部位置的分布，需要约 $2NDV_{\mathrm{voc}}$ FLOP；纯生成的预填充只需最后一个位置的词表分布，可以把这部分降为约 $2DV_{\mathrm{voc}}$，而网络内部仍需计算各层的前缀键值。

式~\eqref{l12:eq:layer-flops} 的二次项来自位置两两匹配与聚合。若实现真正跳过上三角，则 $QK^{\mathsf T}$ 与 $PV$ 的有效计算合计约为 $2N(N+1)D$，常数约减半；直接调用完整的稠密矩阵乘法时，上三角即使在归一化后为零，也不会自动省掉乘法。两种实现仍具有 $\Theta(N^2D)$ 的预填充算术阶数。

模型宽度也决定主导开销。举例而言，若 $N=D$ 且 $D_{\mathrm{ff}}=4D$，完整稠密基线中位置匹配与聚合占主要矩阵乘法算量的 $4/(24+4)=1/7$。长度增加时这一占比上升。因此，判断耗时需要同时考虑 $N$、$D$、前馈宽度和实际核函数效率，不能只根据某一步含有 $N^2$ 就决定优化顺序。

\subsection{大规模投影的参数与容量算例}

以 GPT-3 的大规模配置中的投影尺寸为例，取 $D=12\,288$、$h=96$、$d=128$、$L=96$。这里只统计每层三个注意力输入投影及一个输出投影，忽略偏置。单头的一个输入投影有 $Dd=12\,288\times128=1\,572\,864$ 个参数；全部头的三个投影合计 $3hDd=3D^2$。输出矩阵有 $D^2$ 个参数，因而 $L$ 层共有
\begin{equation}
 n_{\mathrm{proj}}=L(3D^2+D^2)
 =4LD^2
 =57\,982\,058\,496.
 \label{l12:eq:gpt3-params}
\end{equation}
它约为 $58.0\times10^9$ 个注意力投影参数。每个新位置都要使用这些投影，按每次乘加 $2$ FLOP 计，其主要算量为 $2n_{\mathrm{proj}}\approx115.964\times10^9$ FLOP。

\begin{table}[htbp]
\centering
\caption{给定宽度和层数下的投影规模。参数数为精确整数，运算量按乘加计数后取近似显示。}
\label{l12:tab:gpt3}
\small
\begin{tabularx}{\linewidth}{@{}Y L{3.6cm}L{2.6cm}L{3.3cm}@{}}
\toprule
对象 & 参数数 & 16位容量 & 单位置算量\\
\midrule
单头的一种输入投影 & $1\,572\,864$ & $3\ \mathrm{MiB}$ & $3.146\ \mathrm{MFLOP}$\\
单层全部输入投影 & $452\,984\,832$ & $864\ \mathrm{MiB}$ & $0.906\ \mathrm{GFLOP}$\\
单层输出投影 & $150\,994\,944$ & $288\ \mathrm{MiB}$ & $0.302\ \mathrm{GFLOP}$\\
96层四类投影 & $57\,982\,058\,496$ & $108\ \mathrm{GiB}$ & $115.964\ \mathrm{GFLOP}$\\
\bottomrule
\end{tabularx}
\end{table}

容量按每个参数 $2$ 字节计算，其中 $1\ \mathrm{MiB}=2^{20}$ 字节、$1\ \mathrm{GiB}=2^{30}$ 字节；运算量前缀采用十进制，即 $1\ \mathrm{MFLOP}=10^6$ FLOP、$1\ \mathrm{GFLOP}=10^9$ FLOP。显存容量的二进制单位与运算计数的十进制前缀需要分别处理。

这些结果仅来自四类投影矩阵。上下文相关的匹配与聚合、前馈网络、词表投影以及各类中间状态都还未计入，因此 $108\ \mathrm{GiB}$ 不是完整模型的总显存需求，$115.964\ \mathrm{GFLOP}$ 也不是生成一个词元的全部计算量。这个算例的用途是把矩阵尺寸、参数存储和单位置算量对应起来。

\subsection{预填充与解码的资源需求}

预填充中的投影具有许多输入行，同一权重块可以服务多个位置，适合用分块矩阵乘法增加数据复用。单序列增量解码只有一行输入，投影退化为向量与矩阵乘法，单次读入的权重参与的计算更少。批量处理多个序列可以增加同时参与计算的行数，但同时也会增加各序列的缓存容量和键值读取量。

从算术强度（arithmetic intensity）可以看出这个差别。假设一个投影权重矩阵含 $D d$ 个元素，本次需要把它们全部从显存读取，每元素占 $b_{\mathrm{elem}}$ 字节。只处理一个新位置时，主要算量为 $2Dd$，仅对权重读取量计算的算术强度约为 $2/b_{\mathrm{elem}}$ FLOP/字节。若以 $16$ 位保存权重，这个比值约为 $1$ FLOP/字节。多个输入行共同使用权重块时，可以把这一次读取摊到更多乘加操作上。

若同时统计输入读取和输出写入，可进一步量化批量带来的复用。设一个线性层的权重形状为 $m\times n$，处理 $B$ 行输入。假定输入、权重各读取一次，输出写入一次，每个元素占 $b_{\mathrm{elem}}$ 字节，忽略偏置和其他操作，则
\begin{equation}
 AI_{\mathrm{linear}}\approx
 \frac{2Bmn}{b_{\mathrm{elem}}(mn+Bm+Bn)}.
 \label{l12:eq:batch-ai}
\end{equation}
取 $m=1024$、$n=4096$、$b_{\mathrm{elem}}=2$，$B=1$ 时约为 $0.999$ FLOP/字节；$B=512$ 时约为 $315.08$ FLOP/字节。输入行数增多使权重读取被更多计算复用，而输入和输出的数据量也随之增长。因此，这个完整传输估计与前面只统计权重读取的估计具有不同分母。真实实现若重复读取数据，算术强度还会降低。

不同请求可以使用同一组投影权重，但通常必须读取各自的键值缓存。增加请求批量不会自动减少每个查询所需的历史键值读取，也不会自动提高这一部分按缓存读取计算的算术强度；它仍可能通过增加并行工作改善设备利用率。批量扩大时，还需同时检查缓存容量、序列长度差异及调度开销。

注意力读取缓存也具有相似的限制。单头一个新查询的匹配与聚合共约 $4Nd$ FLOP，需要读取约 $2Nd$ 个历史键值元素；只考虑这些读取时，算术强度同样约为 $2/b_{\mathrm{elem}}$。这解释了为什么缓存虽然减少计算，较长上下文的逐词元解码仍可能受显存带宽限制。实际缓存命中、并行批量、数据类型和核函数融合都会影响最终表现。

设某段运算需要 $F$ FLOP、传输 $M$ 字节，设备在相应数据类型上的峰值算术吞吐为 $R_{\mathrm{comp}}$、显存带宽为 $R_{\mathrm{mem}}$，则理想时间至少满足
\begin{equation}
  T\geq\max\!\left(\frac{F}{R_{\mathrm{comp}}},
                   \frac{M}{R_{\mathrm{mem}}}\right).
  \label{l12:eq:resource-bound}
\end{equation}
真实时间还包含无法完全重叠的访存、同步、启动和布局转换等开销。因此，算量减少多少与最终时间减少多少之间，还隔着硬件资源利用和实现方式。

\subsection{分块融合如何减少中间矩阵读写}

基线预填充通常显式生成分数矩阵 $S$ 和权重矩阵 $P$。每个头每张矩阵包含 $N^2$ 个元素，矩阵乘法、归一化和值聚合在相继的核函数中对它们写入和读取。例如 $B=1$、$h=32$、$N=4096$ 时，一张所有头合计的矩阵有 $32\times4096^2=536\,870\,912$ 个元素；若以 $32$ 位保存，仅这一张矩阵就占 $2\ \mathrm{GiB}$，同时保留 $S$ 和 $P$ 则需要 $4\ \mathrm{GiB}$。这部分是单层的临时状态，与跨层持久保存的 KV 缓存承担不同的作用。

闪电注意力（FlashAttention）通过分块与核函数融合，在片上保存局部查询、键值和中间统计量，使完整的 $N\times N$ 分数与概率矩阵不必写回显存。对一组查询，逐块读取键和值，计算局部分数、更新归一化统计量，并累积输出。它改善的是数据搬运和中间存储；完整稠密注意力需要处理的位置关系仍然是二次规模。

分块计算要保持相同的归一化结果，需要在读取后续块时正确调整已经累积的量。对某一个固定查询，设已经处理的位置集合为 $A$，分数为 $s_j$。保存最大值 $m=\max_{j\in A}s_j$、指数和 $\lambda=\sum_{j\in A}e^{s_j-m}$，以及未归一化输出向量 $u=\sum_{j\in A}e^{s_j-m}v_j$。当前输出就是 $u/\lambda$。

假设随后读入另一个非空可见块 $J$，其最大分数为 $m_J$。新的公共最大值取 $m'=\max(m,m_J)$。由于 $e^{s_j-m'}=e^{m-m'}e^{s_j-m}$，旧的累计量只需乘同一个校正因子，就能与新块合并：
\begin{equation}
\begin{aligned}
 m'&=\max(m,m_J),\qquad \alpha=\exp(m-m'),\\
 \lambda'&=\alpha\lambda+\sum_{j\in J}\exp(s_j-m'),\\
 u'&=\alpha u+\sum_{j\in J}\exp(s_j-m')v_j.
\end{aligned}
\label{l12:eq:online-softmax}
\end{equation}
先用第一个非空可见块初始化这些量，再反复使用此递推；因果掩码下没有可见元素的块直接跳过。读完全部可见位置后，输出 $u/\lambda$，即与整行归一化后加权求和相同的数学结果。有限精度下，运算顺序的改变可能带来舍入差异。这个递推说明了为什么只保留行统计量与一个长度为 $d$ 的累加向量，就能避免存储整行的全部概率。

KV 缓存复用不同生成步骤之间已经计算的历史键值，分块融合则减少同一次注意力计算内部的显存往返。这两个方向可以共同使用：预填充高效地建立各层缓存，后续增量解码继续访问这些缓存，并使用适合单查询或小批量形状的核函数。

\subsection{从核函数性能到完整生成时间}

完整推理包括初始分词和嵌入、预填充、重复的增量解码、词表选择以及必要的数据传输。首个词元的延迟主要涉及初始输入处理和预填充；后续逐词元延迟取决于增量路径。测量时应明确输入长度、输出长度、批量和数据类型，并将这两个阶段分别观察，再统计完整生成任务的时间。

局部加速对整体性能的影响由它在总时间中的占比决定。若某部分原先占总时间的比例为 $f$，优化后加速 $r$ 倍，其余部分耗时保持不变，则整体加速比为
\begin{equation}
  R_{\mathrm{total}}=\frac{1}{(1-f)+f/r}.
  \label{l12:eq:amdahl}
\end{equation}
例如某部分占原时间的 $60\%$，即使该部分加速 $4$ 倍，整体加速也只有 $1/(0.4+0.6/4)\approx1.82$ 倍。这个假设算例说明，优化时应结合实际耗时占比评估收益；核函数之间多出的缓冲区转换或同步也可能抵消局部收益。

给定预训练权重的 GPT-2 推理项目可以围绕这条计算链展开。第一阶段把已有 CPU 参考实现对应的基本运算移到 GPU，并建立数值一致性；第二阶段分析单个核函数和完整推理流程的耗时，提出并实现相应优化；第三阶段继续优化并分析完整生成任务的性能。预训练参数保持固定，工作重点是正确实现和加速前向推理。

从计算依赖来看，投影和前馈主要改变特征表示，因果注意力把历史信息汇入当前位置，逐行归一化把分数变为稳定权重，键值缓存让已完成的位置计算得到复用。对应到实现，矩阵乘法的数据复用、归约的同步组织、缓存的按层维护以及中间结果的读写，共同决定完整生成过程的性能。
